\section{Conclusions}\label{sec:conclusions}

This study assembled a reproducible climate-morphology evidence frame for the
Izmir functional urban region and tested what can be learned from global weak
labels plus a completed local audit. A 250 m matrix covering 16,506 grid units
integrates a
local 5 m surface model, exact building and street geometry, vegetation
structure, remote sensing, coast, OSM context, population demand, and
network-cost green-space accessibility. Under five spatial block folds, the
primary 800 m LightGBM reached weak-label macro-F1 0.838, compared with 0.663
for morphology alone. Vegetation produced the largest additive gain, 2SFCA
contributed a modest model-dependent increment, and leave-district-out
performance exposed substantial geographic transfer loss. Weak classes also
showed a large LST contrast and coherent vegetation/built-form associations.

The methodological conclusion is stronger than any single score. LCZ labels are
most useful here as a common vocabulary and sampling scaffold, not as local
ground truth. Fine surface-model information, 3D vegetation, network
accessibility, spatial validation, audit comparison, and explicit proxy language
make that scaffold more planning-relevant and more falsifiable. Euclidean buffers would not represent
the same accessibility construct as 2SFCA, and random splits would not reveal
the district transfer problem.

CLIMORFA therefore provides a reproducible multimodal evidence framework whose
tabular branch is evaluated against spatial folds and a completed manual audit.
The 8.3\% agreement in the primary high-quality four-class population is an
important finding: it identifies class-space compression that should inform
future local subtype work. The architecture, branch comparisons, attribution
tools, calibrated uncertainty, and multi-scale analysis provide a coherent path
for extending the evidence base while retaining the same audit and spatial-
validation principles.
